\begin{table}[!htbp]
\centering\small
\setlength{\tabcolsep}{4pt}
\renewcommand{\arraystretch}{1.12}
\begin{tabularx}{\linewidth}{@{}>{\raggedright\arraybackslash}p{.16\linewidth}>{\raggedright\arraybackslash}X>{\raggedright\arraybackslash}X@{}}
\toprule
Magnitud & Resultado en la realización especificada & Referencia y contraste\\
\midrule
Planck, anterior al retorno
& \(\hbar_{\rm pre}=1.0545718176461544696\ldots\)
  \(\times10^{-34}\,\mathrm{J\,s}\).
& SI: \(\hbar_{\rm SI}=h_{\rm SI}/(2\pi)\), exacto.
  Su coordenada es \(1.0545718176461563912\ldots\)
  \(\times10^{-34}\,\mathrm{J\,s}\).
  Residuo relativo \(-1.822221\times10^{-15}\).\\[4pt]
Planck, posterior al retorno
& \(\hbar_{\rm ret}=1.0545718169150390611\ldots\)
  \(\times10^{-34}\,\mathrm{J\,s}\).
& Misma referencia SI.
  Residuo relativo \(-6.932836\times10^{-10}\).
  Para \(h_a=2\pi\hbar_a\), los residuos son idénticos.\\[4pt]
Barbero--Immirzi
& \(\gamma_{\HMT}=0.2740076726944592747\ldots\).
& Ghosh--Mitra: aproximadamente \(0.2740668582328300\).
  Diferencia relativa \(-2.159529202\times10^{-4}\)
  entre construcciones teóricas.\\[4pt]
Mezcla CKM
& \(\delta_{\rm CP}=65.676173123554^\circ\),
  \(J=3.1189723326632974\times10^{-5}\).
  Los tres ángulos de \eqref{eq:ii-ckm-decimales} son
  \(\theta_{12}=13.003313589509^\circ\),
  \(\theta_{23}=2.398056157332^\circ\) y
  \(\theta_{13}=0.213977382244^\circ\).
& PDG 2025: las cinco comparaciones de
  \(s_{12},s_{23},s_{13},\delta,J\) satisfacen \(|z_x|<0.029\)
  con las incertidumbres marginales de la
  tabla~\ref{tab:ii-ckm-comparacion}.\\[4pt]
Boltzmann
& \(B_*=1.380649\times10^{-23}\,u_E/u_\Theta\),
  con las bases y la regla termométrica de
  \eqref{eq:ii-b-evaluacion-comparada}.
& SI: \(1.380649\times10^{-23}\,\mathrm{J/K}\), exacto.
  Coincidencia del numeral bajo el transporte dimensional declarado.\\[4pt]
Longitud \(L\)
& \(L=1.6162549826251263301\ldots\)
  \(\times10^{-35}\,\mathrm m\), para \(L_*=1\,\mathrm m\).
& Composición longitudinal y radial:
  \(G_{\rm ret}=c_{\rm int}^3L^2/\hbar_{\rm ret}\),
  con las mismas bases de velocidad y acción.\\[4pt]
Gravitación
& \(G_{\rm ret}=6.6742996594140227\ldots\)
  \(\times10^{-11}\,\mathrm{m^3\,kg^{-1}\,s^{-2}}\).
& CODATA 2022: \(6.67430(15)\times10^{-11}\) en las mismas unidades.
  Diferencia de aproximadamente \(-0.00227057\) incertidumbres estándar.\\
\bottomrule
\end{tabularx}
\caption{Síntesis del contraste cuantitativo. Las evaluaciones, las bases
y las referencias son las de las tablas~\ref{tab:ii-planck-comparacion}--\ref{tab:ii-kb-g-cartas}.
Los residuos relativos de las constantes de acción, la comparación
teórica de Barbero--Immirzi y los residuos normalizados de CKM y gravitación
conservan sus respectivos significados.}
\label{tab:ii-conclusiones-contraste}
\par\medskip
\centering\small
\setlength{\tabcolsep}{4pt}
\renewcommand{\arraystretch}{1.18}
\begin{tabular}{@{}lrrr@{}}
\toprule
Observable & Evaluación HMT & PDG 2025 & \(z_x\)\\
\midrule
\(s_{12}\) & \(0.225007404757\) & \(0.22501\pm0.00068\) & \(-0.003817\)\\
\(s_{23}\) & \(0.041841757010\) & \(0.04183^{+0.00079}_{-0.00069}\) & \(0.014882\)\\
\(s_{13}\) & \(0.003734601164\) & \(0.003732^{+0.000090}_{-0.000085}\) & \(0.028902\)\\
\(\delta\) (rad) & \(1.146265461116\) & \(1.147\pm0.026\) & \(-0.028251\)\\
\(10^5J\) & \(3.118972333\) & \(3.12^{+0.13}_{-0.12}\) & \(-0.008564\)\\
\bottomrule
\end{tabular}
\caption{Valores CKM y comparación marginal reunidos en las conclusiones.
Los tres ángulos de la tabla~\ref{tab:ii-conclusiones-contraste} se contrastan
mediante \(s_{ij}=\sin\theta_{ij}\), junto con la fase en radianes y el
invariante \(J\). Se conservan las cifras, incertidumbres y residuos de la
tabla~\ref{tab:ii-ckm-comparacion}, correspondientes al ajuste de tres
generaciones PDG 2025. La última fila expresa conjuntamente valor e
incertidumbre en unidades de \(10^{-5}\). Los residuos son comparaciones
marginales de observables correlacionados.}
\label{tab:ii-conclusiones-ckm}
\end{table}
